\SourceLeaf{535}\par Die bisher von den Rechnern verwendeten Verfahren zur Bestimmung von $\frac{\pi}{4}=(1)$ stellt \textsc{Gauss} in der folgenden Übersicht zusammen
\begingroup\fontsize{10}{13}\selectfont\setlength{\tabcolsep}{3pt}
\begin{longtable}{@{}l p{.74\linewidth}@{}}
\textsc{Machin} & $(1)=4(5)-(239)$\par auch \textsc{Clausen} \\[3pt]
\textsc{Euler} & ${}=(2)+(3)$\par (\textsc{Euler} an \textsc{Goldbach} 1746 Mai 28) \\[3pt]
\textsc{Vega} & ${}=5(7)+2(\frac{79}{3})$\par (\textsc{Vega} Thesaurus logar. p. 633) \\[3pt]
\textsc{Vega} & ${}=2(3)+(7)$\par auch \textsc{Clausen} (Astr. Nachr. Bd. 25. S. 209) \\[3pt]
\textsc{Rutherford} & ${}=4(5)-(70)+(99)$\par (Philos. Trans. 1841. p. 283) \\[3pt]
\textsc{Dase} & ${}=(2)+(5)+(8)$\par (\textsc{Crelle} Journal. Bd. 27. S. 198) \\[3pt]
\textsc{Gauss}. 1. & ${}=12(18)+8(57)-5(239)$ \\[3pt]
\textsc{Gauss}. 2. & ${}=12(38)+20(57)+7(239)+24(268)$ \\[3pt]
\end{longtable}\endgroup
\par Die ersten Rechnungen für die Tafeln stammen aus der Zeit der Ausarbeitung der \textit{Disquiss. Arr.}; sie wurden dann besonders in den Jahren 1846 und 47 vorangebracht. Am 21. Juli 1847 waren 2283 Zerlegungen in der hier wiedergegebenen Ordnung in Tafeln gebracht. Die übrigen 169 wurden später berechnet, und ich habe sie in diesen Abdruck eingeordnet, der sich in der Einrichtung vom Original nur durch die zur Erleichterung des Satzes verwendete Wiederholung der Faktoren anstelle von Potenzen unterscheidet.
\par Die Handschriften mit diesen letzten Rechnungen scheinen die Ergebnisse in der Form zu enthalten, in der sie unmittelbar gefunden wurden. Die Reihenfolge der Zahlen $a$ lässt vermuten, dass nur für die kleineren die Teiler von $aa+1$ usw. aufgesucht wurden und die größeren Zahlen daraus durch besondere Kunstgriffe hervorgingen. Aufgezeichnet ist aber nur die folgende Regel: \textit{Aus drei Zahlen $a$, $2a-n$, $2a+n$ findet sich eine vierte}
\[
\frac{4a^3-(nn-3)a}{nn+1}
\]
\textit{Diese ist für $n=0$ und $n=1$ stets eine ganze Zahl, sonst nur}
\[
\begin{aligned}
&\textit{für}\quad a\equiv0\quad\textit{und}\quad\equiv\pm\surd-1\quad mod(nn+1)\quad\textit{wenn }n\textit{ gerade}\\
&\textit{und für}\quad a\equiv0\quad\textit{und}\quad\equiv\pm\surd-1\quad mod\frac{nn+1}{2}\quad\textit{wenn }n\textit{ ungerade}
\end{aligned}
\]
\begin{center}
\begin{tabular}{lrrr}
\textit{Beispiele} & $a=253,$ & $n=6,$ & $1750507$\\
 & $a=294,$ & $n=11,$ & $832902$\\
 & $a=119,$ & $n=1,$ & $3370437$\\
 & $a=57,$ & $n=3,$ & $74043$\\
 & $a=123,$ & $n=9,$ & $90657$
\end{tabular}
\end{center}
\par Zur vierten Zahl gehören nämlich keine anderen Primteiler als zu den ersten drei; davon sind nur diejenigen ungeraden Primteiler ausgeschlossen, die zur Zahl $n$ gehören.
\par Die Tafel ‘\textit{Quadratorum myrias prima}’ enthält in der Überschriftenzeile die Tausender und Hunderter, in der ersten senkrechten Spalte die Zehner und Einer der Grundzahl, ferner in der letzten
